%==============================================================================
%  ملحق - الأمان والمتانة
%==============================================================================

%==============================================================================
\chapter{ملحق --- الأمان والمتانة}
\label{chap:securite}
%==============================================================================
الأمان متطلب \textbf{مُستعرِض} موجود سلفًا في كراس الشروط: تشفير التبادلات بـ
\textbf{X.509 / SSL}، ومصادقة مُفوَّضة بـ \textbf{OAuth Google}، ومصفوفة حقوق
\textbf{RBAC}، والامتثال لـ \textbf{RGPD} (الملحق~\ref{chap:complements}).
يُوحّد هذا الملحق ذلك في استراتيجية متّسقة لـ\textbf{الدفاع في العمق}، ثم يعالج
\textbf{المتانة} (مقاومة الحِمل والذُّرى)، لأن \emph{التوفّر} إحدى خصائص الأمان
الثلاث. تتّبع الممارسات الجيدة المعتمَدة أدلّة المجال المرجعية (تصليب موقع ويب)
والأداة الحرة \textbf{k6} لاختبارات الحِمل.

\begin{encadre}
\textbf{ثالوث CIA:} يهدف الأمان إلى \textbf{السرّية} (التشفير، التحكّم في
الوصول)، و\textbf{السلامة} (التحقّق من المُدخلات، الاستعلامات المُحضَّرة، النسخ
الاحتياطي)، و\textbf{التوفّر} (مكافحة DDoS، تحديد المعدّل، اختبارات الحِمل). لا
لبنة منها خدمة خارجية \emph{إجبارية}: لكلٍّ مكافئ حرّ وقابل للنشر الذاتي
(ModSecurity، Let's Encrypt، k6)، وفق متطلب النشر الذاتي.
\end{encadre}

%------------------------------------------------------------------------------
\section{نموذج تهديد تركيبي}
%------------------------------------------------------------------------------
يضع الجدول المخاطر المعتادة لتطبيق ويب في مقابل الإجراء المعتمَد في Zümm.
\begin{xltabular}{\textwidth}{>{\bfseries}p{4cm} X}
\toprule
\textbf{التهديد} & \textbf{الإجراء المعتمَد في Zümm} \\
\midrule
\endhead
اعتراض التبادلات & TLS / SSL من الطرف إلى الطرف، شهادات X.509، إعادة توجيه HTTPS وترويسة HSTS. \\
سرقة بيانات الاعتماد & OAuth Google (لا كلمة سرّ مُخزَّنة)، مصادقة ثنائية يحملها المزوّد. \\
حقن SQL & استعلامات مُحضَّرة (\emph{PreparedStatement}) عبر JDBC، تحقّق وتنميط المُدخلات. \\
XSS (سكربت محقون) & تهريب منهجي للمُخرَجات، ترويسة \texttt{Content-Security-Policy}. \\
CSRF (طلب مزوَّر) & رمز مضادّ لـ CSRF لكل جلسة، كعكات \texttt{SameSite}. \\
اختطاف الجلسة & كعكات \texttt{HttpOnly} و\texttt{Secure}، انتهاء الجلسة وتدويرها (سلفًا §\,7). \\
تصعيد الامتيازات & تحكّم في الوصول \textbf{RBAC} وأدنى امتياز (الملحق~\ref{chap:complements}). \\
حجب الخدمة (DDoS) & تحديد المعدّل، جدار حماية تطبيقي (WAF)، استيعاب لا تزامني للقياسات. \\
كشف المعطيات & تشفير عند السكون، تقليل RGPD، نسخ احتياطي مشفّر. \\
\bottomrule
\end{xltabular}

%------------------------------------------------------------------------------
\section{الدفاع في العمق}
%------------------------------------------------------------------------------
لا يعتمد الأمان على حاجز وحيد بل على \textbf{طبقات متتالية}: إذا سقطت واحدة،
صمدت الأخرى. يُسقِط الشكل~\ref{fig:secu-couches} هذه الطبقات على بنية Zümm.

\begin{figure}[htbp]
\centering
\begin{tikzpicture}[
  couche/.style={draw=mielfonce,thick,rounded corners=3pt,minimum width=12.2cm,minimum height=8mm,font=\footnotesize,align=center},
  l1/.style={couche,fill=grisdoux},
  l2/.style={couche,fill=miel!12},
  l3/.style={couche,fill=miel!22},
  l4/.style={couche,fill=miel!34},
  l5/.style={couche,fill=miel!12}
]
\node[l1] (c1) at (0,4) {\textbf{المحيط}: WAF، مكافحة DDoS، تحديد المعدّل، CDN (اختياري، قابل للنشر)};
\node[l2] (c2) at (0,3) {\textbf{النقل}: TLS / SSL، شهادات X.509، إعادة توجيه HTTPS، HSTS};
\node[l3] (c3) at (0,2) {\textbf{التطبيق}: OAuth2، RBAC، التحقّق من المُدخلات، ترويسات CSP، مضادّ CSRF};
\node[l4] (c4) at (0,1) {\textbf{المعطيات}: استعلامات مُحضَّرة، تشفير عند السكون، أدنى امتياز لقاعدة المعطيات};
\node[l5] (c5) at (0,0) {\textbf{الاستغلال}: التحديثات، التسجيل والإشراف، نسخ احتياطي مُختبَر};
\end{tikzpicture}
\caption{الدفاع في العمق: خمس طبقات حماية مُسقَطة على بنية Zümm.}
\label{fig:secu-couches}
\end{figure}

\subsection{المحيط والنقل}
على الحافة، يُرشّح \textbf{جدار حماية تطبيقي} (WAF، مثلًا ModSecurity) الطلبات
الخبيثة، ويُخمِّد \textbf{تحديد المعدّل} الذُّرى ومحاولات الإساءة. تشفير
\textbf{النقل} مطلوب سلفًا: TLS / SSL، شهادات \textbf{X.509}، إعادة توجيه إجبارية
إلى HTTPS وترويسة \texttt{Strict-Transport-Security} (HSTS) لمنع أي تراجع إلى
نصّ صريح.

\subsection{التطبيق}
تجمع طبقة التطبيق \textbf{المصادقة} المُفوَّضة (OAuth Google، بعامل ثانٍ يحمله
المزوّد)، و\textbf{التحكّم في الوصول} RBAC بأدنى امتياز
(الملحق~\ref{chap:complements})، و\textbf{التحقّق من المُدخلات}، و\textbf{الحماية
المضادّة لـ CSRF}، و\textbf{ترويسات الأمان} (\texttt{Content-Security-Policy}،
\texttt{X-Frame-Options}، \texttt{X-Content-Type-Options}). تعتمد إدارة الجلسة
على كعكات \texttt{HttpOnly} و\texttt{Secure} المُحدَّدة سلفًا في
الفصل~\ref{chap:conception}.

\subsection{المعطيات والاستغلال}
يمرّ الوصول إلى المعطيات حصريًا عبر \textbf{استعلامات مُحضَّرة} (مضادّة للحقن)،
وحساب قاعدة \textbf{بأدنى امتياز}، و\textbf{تشفير عند السكون} للمعطيات الحسّاسة.
في الاستغلال، تُغلِق \textbf{التحديثات} المنتظمة للخادم والتبعيات،
و\textbf{التسجيل} و\textbf{الإشراف}، إضافةً إلى \textbf{نسخ احتياطي مشفّر تُختبَر
استعادته}، الحلقة.

%------------------------------------------------------------------------------
\section{شبكة التصليب}
%------------------------------------------------------------------------------
يُحصي الجدول الممارسات الجيدة لتصليب موقع ويب ويحدّد، لكلٍّ، \textbf{حالتها} في
Zümm: مطلوبة سلفًا من كراس الشروط أو مُضافة بهذا الملحق.
\begin{xltabular}{\textwidth}{>{\bfseries}p{4.2cm} X C{2.2cm}}
\toprule
\textbf{الممارسة} & \textbf{التنفيذ في Zümm} & \textbf{الحالة} \\
\midrule
\endhead
HTTPS / TLS في كل مكان & شهادات X.509، إعادة توجيه HTTPS، HSTS. & مطلوب سلفًا \\
مصادقة قوية & OAuth Google + عامل ثانٍ من المزوّد. & مطلوب سلفًا \\
التحكّم في الوصول (RBAC) & أدوار وأذونات بأدنى امتياز. & مطلوب سلفًا \\
استعلامات مُحضَّرة & وصول JDBC مُعامَل حصريًا. & مطلوب سلفًا \\
جدار حماية تطبيقي (WAF) & ترشيح الطلبات على المحيط (ModSecurity). & \textbf{مُضاف} \\
تحديد المعدّل / مكافحة DDoS & تخميد الذُّرى والإساءة. & \textbf{مُضاف} \\
ترويسات الأمان & CSP، X-Frame-Options، X-Content-Type-Options. & \textbf{مُضاف} \\
حماية مضادّة لـ CSRF & رمز لكل جلسة، كعكات SameSite. & \textbf{مُضاف} \\
التحديثات والتصحيحات & خادم التطبيقات والتبعيات مُحدَّثة. & \textbf{مُضاف} \\
التسجيل والإشراف & آثار الوصول والأخطاء، تنبيهات. & \textbf{مُضاف} \\
نسخ احتياطي مشفّر & نسخ منتظم + استعادة مُختبَرة. & \textbf{مُضاف} \\
تحليل مضادّ للبرمجيات الخبيثة & فحص دوري للمستودعات والتحميلات. & \textbf{مُضاف} \\
\bottomrule
\end{xltabular}

%------------------------------------------------------------------------------
\section{المتانة: اختبارات الحِمل والموثوقية (k6)}
%------------------------------------------------------------------------------
لمّا كان \textbf{التوفّر} خاصية أمان، وجب أن يقاوم النظام الحِمل الاسمي
و\textbf{الذُّرى}. نعتمد \textbf{k6}، أداة اختبار حِمل \textbf{حرة}، قابلة
للبرمجة، تُحاكي \emph{مستخدمين افتراضيين} (VUs) وتتحقّق من \textbf{عتبات}
موضوعية (\emph{thresholds}).

\subsection{أنواع الاختبارات}
\begin{xltabular}{\textwidth}{>{\bfseries}p{3cm} X}
\toprule
\textbf{النوع} & \textbf{الهدف} \\
\midrule
\endhead
الحِمل (\emph{load}) & التحقّق من السلوك عند الحِمل المتوقّع (مستخدمون وقياسات اسمية). \\
الإجهاد (stress) & إيجاد نقطة الانهيار بزيادة الحِمل فوق الاسمي. \\
الذروة (\emph{spike}) & محاكاة دفقة مفاجئة، مثلًا تدفّق ضخم لقياسات المستشعرات. \\
التحمّل (\emph{soak}) & كشف تسرّبات الذاكرة والتدهور على مدة طويلة. \\
\bottomrule
\end{xltabular}

\subsection{سيناريوهات خاصّة بـ Zümm}
\begin{itemize}
  \item \textbf{استيعاب بالدفقات}: \texttt{POST /api/mesures} بالحِزم (شقّ المستشعرات، §\,4.5)، للتحقّق من الطابور اللا تزامني وعدم التكرار.
  \item \textbf{لوحات القيادة}: استعراض متزامن للوحتي الإنتاج والتنبيهات (تجميعات، رسوم).
  \item \textbf{خدمة ويب API}: استدعاءات متكرّرة لـ \texttt{getZummHoneyActualQuantity} من تطبيقات خارجية.
\end{itemize}

\subsection{العتبات والحكم}
يفشل الاختبار إذا لم تُحقَّق العتبات الموضوعية، ما يدلّ على عنق زجاجة يجب تصحيحه
(حجم المجمع، الفهرس، الذاكرة المؤقتة، التوسيع). للتوضيح:
\begin{verbatim}
thresholds:
  request p95 latency        < 500 ms
  HTTP error rate            < 1 %
  /api/mesures ingestion     sustained nominal throughput
\end{verbatim}

\begin{figure}[htbp]
\centering
\includegraphics[width=\linewidth]{diagramme_sequence_test_charge}
\caption{اختبار حِمل k6: تصاعد تدريجي للحِمل، قياس زمن الاستجابة ومعدّل الأخطاء،
والحكم في ضوء العتبات.}
\label{fig:secu-k6}
\end{figure}

\subsection{تكملة لخطة الاختبارات}
تُمدّد الحالات التالية خطة اختبارات الفصل~\ref{chap:tests} (TC01--TC16) على
جانب الأمان والمتانة.
\begin{xltabular}{\textwidth}{>{\bfseries}p{1.6cm} p{4.4cm} X}
\toprule
\textbf{المعرّف} & \textbf{الموضوع} & \textbf{النتيجة المنتظرة} \\
\midrule
\endhead
TC17 & اختبار حِمل (k6) & احترام عتبتي p95 ومعدّل الأخطاء عند الحِمل الاسمي. \\
TC18 & اختبار ذروة (spike) & يمتصّ النظام دفقة دون فقدان أو خطأ دائم. \\
TC19 & حقن SQL / XSS & رفض المُدخلات الخبيثة أو تحييدها (تهريب). \\
TC20 & نسخ احتياطي / استعادة & استعادة كاملة مُتحقَّق منها من نسخة احتياطية مشفّرة. \\
\bottomrule
\end{xltabular}

%------------------------------------------------------------------------------
\section{قائمة تحقّق ما قبل النشر (go-live)}
%------------------------------------------------------------------------------
قبل أي إطلاق في الإنتاج، تُتحقَّق النقاط أدناه. تُلخّص هذه القائمة الأقسام
السابقة في مراجعة تشغيلية، مجالًا مجالًا.
\begin{xltabular}{\textwidth}{>{\bfseries}p{3.4cm} X}
\toprule
\textbf{المجال} & \textbf{نقاط للتحقّق قبل الإنتاج} \\
\midrule
\endhead
الأسرار والإعداد & لا سرّ بنصّ صريح في المستودع؛ \texttt{ConfigZumm.ini} والمفاتيح والشهادات خارج التحكّم في الإصدارات، مُخزَّنة بشكل محميّ. \\
النقل & HTTPS إجباري، إعادة توجيه الحركة الصريحة، ترويسة HSTS نشطة، شهادات X.509 صالحة وغير منتهية. \\
المصادقة والوصول & OAuth Google فاعل، مصفوفة RBAC مُتحقَّقة، الحسابات والوصول الافتراضية مُعطَّلة. \\
التطبيق & ترويسات الأمان (CSP، X-Frame-Options)، حماية مضادّة لـ CSRF، تحقّق من المُدخلات، رسائل خطأ دون تسريب معلومات. \\
المعطيات & استعلامات مُحضَّرة في كل مكان، تشفير عند السكون للمعطيات الحسّاسة، نسخ احتياطي حديث \textbf{واستعادة مُختبَرة}. \\
التبعيات & خادم التطبيقات والمكتبات مُحدَّثة، فحص ثغرات التبعيات مُنجَز. \\
المتانة & اختبار حِمل k6 \textbf{أخضر} (عتبتا p95 ومعدّل الأخطاء)، اختبار ذروة ناجح، خطة تصاعد وتراجع (rollback) مُحدَّدة. \\
القابلية للمراقبة & تسجيل الوصول والأخطاء، إشراف وتنبيهات قائمة، توقيت زمني متّسق. \\
الامتثال & احترام RGPD (التقليل، مدة الحفظ، الحقوق)، البيانات والسجلّ مُحدَّثة (الملحق~\ref{chap:complements}). \\
\bottomrule
\end{xltabular}

\begin{encadre}
\textbf{قاعدة الإطلاق في الإنتاج:} نقطة غير مُحقَّقة \textbf{مانعة} حتى التصحيح.
تُعاد القائمة عند كل تسليم كبير، بروح التحسين المستمرّ وتحت \textbf{مراقبة بشرية}.
\end{encadre}

%------------------------------------------------------------------------------
\section{الامتثال والاتّساق مع المتطلبات}
%------------------------------------------------------------------------------
\begin{encadre}
\textbf{الخلاصة:} يُوحّد الملحق أمانًا \textbf{بالطبقات} (المحيط، النقل،
التطبيق، المعطيات، الاستغلال) ويضيف \textbf{متانة مقيسة} باختبارات الحِمل.
يُمدّد، دون تناقض، المتطلبات المطروحة سلفًا (X.509 / SSL، OAuth، RBAC، RGPD)
ويحترم \textbf{النشر الذاتي}: كل لبنة معتمَدة (WAF من ModSecurity، شهادات Let's
Encrypt، k6) \textbf{حرة ومجانية}. يبقى الأمان، كالتنبيهات، مسارًا مستمرًّا تحت
\textbf{مراقبة بشرية}.
\end{encadre}
